\subsection{The Hypotheses and Design Decisions Behind the Search}
\label{sec:hypotheses}

The exploratory analysis of Sections~\ref{sec:eda-ns} and~\ref{sec:eda-bh} fixes what the data allow: how closely a smooth interpolation between neighbouring curves reproduces a curve held out of them. That bound turns the design of a surrogate into one question per dataset and target. \emph{Which surrogate comes within one significant figure of the reachable relative error on the interior curves, holds on the edge curves it was fitted without, keeps the slope of the data, and at what cost to fit and to call?} The question is asked of the two maps of this work -- black holes $(\beta, \rh) \to (M, \Dch)$ and neutron stars $(\rhoc, \beta, \lambda) \to (M, \Dch)$ -- with one model per target, the mass $M$ and the charge target $Y = \log_{10}(\Dch/M)$.

\paragraph{The search space has five axes.}
They are fixed in this order, because the earlier ones carry across models and the later ones depend on the model (the pipeline above):
\begin{enumerate}
  \item \emph{representation}: the scale of the inputs, the target of each model ($M$ or $\log_{10} M$; $Y$ or $\Dch$ itself), and whether a row or a whole curve is the unit of prediction (Section~\ref{sec:representation});
  \item \emph{model family}: lookup and interpolation between curves, Gaussian process regression, gradient-boosted trees, multilayer perceptrons and residual networks (Section~\ref{sec:families});
  \item \emph{hyperparameters} of each surviving family, each search stopped at a time budget of $\knobTuningBudgetMinutes$ minutes and given a verdict: searched, grid when fixed points are scored in place of a search, or cut when the search ended on the clock (Section~\ref{sec:precision});
  \item \emph{training regime}: the loss, the activation, the optimizer and the numerical precision (Section~\ref{sec:precision});
  \item \emph{data budget}: how the error falls with the number of training curves.
\end{enumerate}
A candidate is scored by the significant figures it keeps on held-out curves.

\paragraph{Three choices the analysis suggests are tested on their own.}
\begin{itemize}
  \item \textbf{$H_1$, the log target (representation).} Regressing $Y$ instead of $\Dch$ makes the squared error proportional to the squared relative error of the charge rebuilt with the true mass,
  \[
    \left| Y_{\mathrm{pred}} - Y_{\mathrm{true}} \right|^2 \approx \left( \frac{\Dch_{\mathrm{pred}} - \Dch_{\mathrm{true}}}{\Dch_{\mathrm{true}} \ln 10} \right)^2 ,
  \]
  so one model can keep its relative error flat across the $\nsEdaDOrders$ orders of magnitude the neutron-star charge spans, without a classify-then-regress pipeline. Section~\ref{sec:representation} tests it against regressing $\Dch$ directly, per order of magnitude of $\Dch$. It held on the whole charge range: the log target serves every model, where the linear charge drives the network below zero figures; the split per order of magnitude follows with the families.
  \item \textbf{$H_2$, the turning point (training regime).} Combining the Huber loss ($\delta = 1.0$)~\citep{Huber1964} with smooth GELU activations~\citep{Hendrycks2016_GELU} prevents gradient saturation at the neutron-star turning point $dM/d\rhoc = 0$ and halves the mass error near $\Mmax$ against the mean squared error with ReLU. Section~\ref{sec:precision} does not test it, since no pair winner is a network; the black holes have no interior maximum of the mass, so $H_2$ concerns the neutron stars alone.
  \item \textbf{$H_3$, smoothing across the grid (model family).} Identity skip connections, $h^{(l+1)} = h^{(l)} + F(h^{(l)})$~\citep{He2016_ResNet}, let a network follow the continuous variation along a curve while interpolating smoothly between the discrete grid of curve parameters, $(\beta, \lambda)$ for neutron stars and $\beta$ for black holes, without ripples at the grid steps. Section~\ref{sec:families} tests it with a ripple measure against a plain multilayer perceptron of matched size. It stays open: the screen never kept the perceptron and the residual network on the same pair, so their ripples do not compare.
\end{itemize}

\paragraph{Every tuned candidate must pass a set of gates.}
The four gates are the physics gate, which counts the reversed steps of the test curves, those whose predicted slope has another sign than the slope of the data; the interior gate, the interior figures against the ceiling of the pair; the edge gate, which sets the margin of a map from the curves held out of a refit; and the cost gate, the time budget of $\knobTuningBudgetMinutes$ minutes a search, where a candidate too slow is put on fixed points or cut and never dropped unreported (Sections~\ref{sec:error-metric} and~\ref{sec:validity}).
